% from main.tex
%%%%%%%%%%  scan idx 625  |  folio 612  %%%%%%%%%%
Il
existe un monomorphisme canonique de \(\mathrm{Deff}\) dans le foncteur \(\mathrm{Rec}\) des données de recollement sur \(f^*L\). Par définition \(\mathrm{Rec}\) est égal au foncteur \(\Isom(p_1^*L,p_2^*L)\) où les \(p_i:X\times_YX\to X\) sont les projections, qui est représenté par un schéma affine \(I\) de présentation finie sur \(S\) d'après 2.4. Il reste donc à représenter le monomorphisme
\[
  i:\mathrm{Deff}\hookrightarrow \Isom(p_1^*L,p_2^*L)
\]
par un morphisme de présentation finie, qui sera nécessairement quasi-affine (d'après EGA IV 8.11.2).

Remarquons en passant que si l'on y tient on peut restreindre \(\Isom(p_1^*L,p_2^*L)\) en imposant la condition de cocycle, mais ce n'est pas indispensable pour la suite.

Après le changement de base \(I\to S\), on est donc ramené au problème suivant : si \(L\) est un faisceau inversible sur \(X\) muni d'une donnée de recollement \(\varphi:p_1^*L\xrightarrow{\sim}p_2^*L\), représenter le sous-foncteur de \(S\) «~qui rend cette donnée effective~». Mais on a le

\textbf{Lemme 4.3.} Sous les conditions de 4.1, soit \(L\) un faisceau inversible sur \(X\) muni d'une donnée de recollement
\[
  \varphi:p_1^*L\xrightarrow{\sim}p_2^*L.
\]
Pour que cette donnée soit effective, il faut et il suffit que les cinq conditions suivantes soient réalisées.
\begin{enumerate}[label=\arabic*)]
\item Pour tout \(i\ge0\), \(R^if_*(L)\) est plat sur \(S\).
\item Pour tout \(j\ge0\), \(R^jg_*(p_1^*L)\) est plat sur \(S\).
\item Le conoyau de \(\pi_1-\pi_2:f_*(L)\longrightarrow g_*(p_1^*L)\simeq g_*(p_2^*(L))\) est plat sur \(S\), (ou le dernier isomorphisme est \(g_*(\varphi)\)).
\item Le noyau \(N\) de \(\pi_1-\pi_2:f_*(L)\to g_*(p_1^*L)\) est inversible sur \(Y\).
\item Le morphisme composé \(f^*N\to f^*f_*L\to L\) est un isomorphisme.
\end{enumerate}
